%% notes-tex/wet-lab/env-case-study-gate/sections/2-mockups.tex
%% SOURCE: the sheets are written by experiments/040-inhibitor-synergy-wetlab/scripts/
%% figure_gate_panels.py from the result files named in the source table under each
%% figure; every number in a caption is a macro of tables/gate_numbers.tex written by the
%% same script. Result files are under
%% experiments/040-inhibitor-synergy-wetlab/results/ unless a path is given; 033/, 035/
%% and 038/ abbreviate the sibling experiment folders, which live on unlanded branches
%% and are read by absolute path (the paths are in results/figure_gate_manifest.json).

\clearpage
\section{The mockups\secstatus{todo}}
\label{sec:mockups}

Each caption gives, per panel, what is plotted, the statistic and the status; the table
under each figure gives the files behind it. The statistic on the public side is always
the per-compound centered Spearman under compound-cold folds, median over compounds,
with \hydnFitCompounds{} compounds fitted per fold; on the wet-lab side it is the growth
AUROC over the \hydnCombinations{} combinations and the fitness Spearman over the
\hydnGrewServed{} that grew under the served call. \colorkey{}

\clearpage
\tcfig{\figConcept}{\textbf{The task, the data and the evaluation.} Three
blocks, each a layout rather than a drawn result; a dashed box is proposed and not run.
\textbf{The task} \status{todo}\,\fstat{idea}: a compound and its molar dose enter the
cell graph transformer as one token placed in front of the gene tokens, which attend to
it; a public-screen strain is read at its deleted genes, giving a deletion profile, and
the bAID host, which carries no deletion, is read at the mean over genes, giving growth
and fitness. A mixture is a set of compound tokens, one per inhibitor: each is read
alone through the shared head and the readings compose as Bliss independence (rounds
nine and ten), which is what took the private readout from chance to the rule's level.
\textbf{The data} \status{todo}\,\fstat{file}: the pooled chemogenomic table over
\hydnScreens{} screens, \hydnPoolCells{} cells and \hydnPoolCompounds{} compounds across
\hydnPoolGenes{} genes, of which \hydnPoolCellsWithMolar{} cells carry a molar dose,
beside the private Bioscreen runs, \hydnWells{} wells over \hydnInhibitors{} inhibitors
on one strain. \textbf{The evaluation} \status{todo}\,\fstat{idea}: compound-cold folds
and the per-compound centered Spearman on the public side, growth and fitness against
Bliss and Loewe on the wet-lab side, and the industrial reading, which is to rank the
worst combination scenarios of a hydrolysate.}{fig:concept}

%% The float is flushed before its source table is set: a deferred float
%% otherwise lands on the page AFTER the table that describes it.
\clearpage

{\footnotesize\noindent
\begin{tabular}{@{}p{22mm}p{154mm}@{}}
\toprule
\multicolumn{2}{@{}l}{\textbf{Sources for \cref{fig:concept}}} \\
\midrule
the task & Nothing drawn from a file. The encoder it describes is the environment-encoder
head of \file{038/}; the host head is the proposed adaptation of
\file{experiments/040-inhibitor-synergy-wetlab/scripts/train_mixture.py}. Status: idea. \\
the data & \file{033/pool_totals.csv}, \file{033/store_axes.csv} (the screens, the
cells, the compounds, the genes, the molar-dose coverage); \file{wetlab_wells.csv} and
\file{wetlab_provenance.json} (the wells, the runs, the raw-mirror shas). Status: on
file. \\
the evaluation & \file{038/ladder/ladder_r2_scores.csv} (the fold size);
\file{mixture_scores.csv} (the two bars). The boxes themselves are a layout, so the
block is marked idea even though its counts are on file. \\
\bottomrule
\end{tabular}}

\clearpage
%% TRUE SIZE (\tcfig, no scaling): the sheet is authored at Nature's full width and
%% its height cap, 179 by 170 mm, so it sits here exactly as it would print. The
%% caption carries the sheet's points and one sentence per panel (the standard of
%% 2026-10-18: the title states the takeaway, the body the main points, the secondary
%% points, then the detail); the paragraph after the figure holds the panel-by-panel
%% detail with each panel's status, and the table after that the files behind it.
\tcfig{\figMain}{\textbf{Composition is what the wet lab buys, and no public screen
lifts the public score.} The six hydrolysate inhibitors are met by one wild-type host;
Vanacloig 2022 profiles two of them itself, and the nearest profiled compound of each of
the other four sits beside it (\textbf{a}). Four public deletion screens of three
genotypes and this study's wild-type wells differ in size, culture and in what a label
is, a strain-level score on a screen's own scale against a growth rate per well
(\textbf{b}); what joins them is one gene table, one compound table and one growth
reading, with a screen's IC30 dose asserting the wild type at \hydicThirtyFitness{}
(\textbf{c}). The cell graph transformer reads one token per inhibitor through one head,
the host as the empty genotype at the gene mean, the dose as a gain on the medium's
offset, and composes the readings as Bliss independence (\textbf{d}).
On the public side nested ridge on fingerprint counts, \hydladderRidge{}, is not beaten by
any transformer arm (\textbf{e}, \textbf{f}), ridge gains about \hydcurvePerDoubling{} per
doubling of fitted compounds with the environment encoder below it at every count
(\textbf{g}), and adding screens leaves the compound-cold score flat (\textbf{o}, left). On
the wet-lab side Loewe additivity from the titrations orders the \hydnCombinations{}
combinations at a growth AUROC of \hydloeweAuroc{} and Bliss from the measured singles at
\hydblissAuroc{}; the model trained on the anchors and the titration reads \hydmodelAuroc{}
on combinations it never saw and its front sits inside the isobole grids (\textbf{m},
\textbf{n}), and the zero-shot arms, which never see a titration well, rise with every
screen added, inside the run spread (\textbf{o}, right). Three readings qualify the join.
The host needs \hydratioFf{} and \hydratioHmf{} times the screen's IC30 of furfural and
5-HMF (\textbf{h}); a screen's label is a strain's growth against its own control on its
own scale, this study's a wild-type growth rate with no growth in a share of
\hydlabelNoGrowthHostexTwoThree{} of the combination wells (\textbf{b}); and the public
profiles do not reach the private question: pairs compose as the mean of their singles
(\textbf{i}), furfural's profile is predicted at \hydprofFfRidge{} where its screen
repeats at \hydprofFfReliability{} while 5-HMF's screen does not repeat (\textbf{j}), and
the similarity of two inhibitors' profiles does not anticipate their pair's departure from
Bliss, Spearman \hydsvdBestSpearman{} at p \hydsvdBestP{} (\textbf{k}). The titrations
themselves (\textbf{l}) are read by the model only with the dose offset, Spearman
\hydexTwoOneDoseSpearman{} against \hydexTwoOneFittedSpearman{} without it. Color is the
source of a value (key above the sheet); dark and light fills in \textbf{b} and \textbf{c}
are trained on and held out. Public scores are compound-cold centered Spearman, median over
held-out compounds; wet-lab scores are the served growth call.}{fig:main}
%% the sheet and its caption take a page of their own, before the panel paragraph
\clearpage

\textbf{Panels of \cref{fig:main}.}
\textbf{a}, The six inhibitors from their curated SMILES at the ACS 1996 drawing
standard at half the standard's bond length (0.19 cm; a 0.6 pt line, double bonds
spaced at 18 percent of the bond, Arial labels, implicit hydrogens on the heteroatoms
only; RDKit's ACS mode through \texttt{torchcell.utils.acs\_structure}). The top row,
on the measured-value field, is the two inhibitors Vanacloig 2022 profiles itself,
furfural and 5-HMF; each of the other four stands beside the profiled compound nearest
it by FCFP4 count Tanimoto over the pool's curated SMILES, that column on the field with
the similarity after the name: sodium acetate to acetic acid
(\hydneighborAaSimilarity{}) and to formic acid (\hydneighborFaSimilarity{}), azelaic
acid to levulinic acid (\hydneighborLvaSimilarity{}), ethanol to lactic acid
(\hydneighborLaSimilarity{}). The hydrolysate's two families are the furan aldehydes
from sugar dehydration (pentoses to furfural, hexoses to 5-HMF) and the weak acids
(5-HMF opens to levulinic and formic acid, acetyl groups release acetic acid). After
each inhibitor's name, the host's half-maximal dose from its titration fit.
\status{todo}\,\fstat{file}
\textbf{b}, Left, what the mixture trainer reads from a four-source run, one bar per
source on a log scale: a screen's cells, this study's wells, the anchors' conditions,
with the genotype on each row's second line (Vanacloig's haploid deletions on the
PDR1, PDR3 and SNQ2 sensitized host; Hoepfner's homozygous and Hillenmeyer's
heterozygous deletions in diploids; Wildenhain's haploid deletions; this study's
wild-type BY4742-iAID6 host with no deletion) and the compound count and the culture
beside each bar. The task's profiles are
\hydnVanacloigCells{} anaerobic, static cells over \hydnVanacloigCompounds{} compounds;
the auxiliary screens add \hydnHoepfnerCells{}, \hydnHillenmeyerCells{} and
\hydnWildenhainCells{} aerobic cells, \hydnFourSourceCells{} in all. The host readout is
aerobic, shaken YPD, a different oxygen regime from the profiles. The anchors are the
\hydnAnchorsFour{} public conditions whose served dose is an IC30 (\hydnAnchorCompounds{}
compounds of Vanacloig and Hoepfner, and the compound-free medium), at an asserted, not
measured, fitness. Right, what a label is. A screen's row is a histogram of the value
the trainer reads before it standardizes the source: one deletion strain's growth under
one compound at one dose, relative to that screen's control, oriented so that a sick
deletion is negative, on the screen's own scale (Vanacloig a log2 ratio to the
inhibitor-free control, Hoepfner a MAD-scaled sensitivity score, Hillenmeyer a negated
fitness-defect log2 ratio, Wildenhain a z-score in IQR units), the band spanning its
0.5th to 99.5th percentile with zero marked. This study's rows draw the labels as they
are, on one fitness scale of 0 to 1.6 with 1 ticked: the six titration curves with their
wells, the \hydnCombinations{} combinations as bars sorted by fitness with no growth at
zero, and the three isobole grids as cells; no growth in a share of
\hydlabelNoGrowthHostexTwoOne{} of the titration wells,
\hydlabelNoGrowthHostexTwoThree{} of the combination wells and
\hydlabelNoGrowthHostisobole{} of the isobole wells. The anchors' row is a sketch of
what is asserted and never measured: fitness 1 with no compound and
\hydicThirtyFitness{} at a screen's IC30 dose. The anchors and the titration are
trained on; the combinations and the isobole grids never. \status{todo}\,\fstat{file}
\textbf{c}, What the sources share and what is trained on and tested on, as three
drawings. Genes: every screen is a deletion collection of the same genes, so one table
of \hydnGeneTokens{} gene tokens serves all of them, each screen's coverage drawn as a
bar of that length (\hydjoinGenesMin{} to \hydjoinGenesMax{}); this study's host
deletes none. Compounds: every compound is one row of one structure table (FCFP4
counts), screens share \hydjoinSharedMin{} to \hydjoinSharedMax{} compounds, and of
this study's six inhibitors the \hydjoinInhibitorsInScreens{} filled squares (furfural
and 5-HMF) are in Vanacloig, the other four in no screen. Growth: every value is a
growth reading, a screen's of a strain against its control and this study's of the
wild type, and an IC30 dose asserts the wild type at \hydicThirtyFitness{}, the marked
point of the sketch, which is what ties a screen's dose scale to the host's. The split,
in the fills of b: the trainer sees every screen, with Vanacloig in compound-cold folds
(one of five held out, \hydnVanacloigCompounds{} compounds), the \hydnAnchorsFour{}
anchors and the titration (\hydnTitrationConditions{} conditions,
\hydnTitrationWells{} wells); it is tested on the held-out compounds' profiles
(centered Spearman over genes) and on this study's \hydnCombinations{} combinations
(growth AUROC, fitness Spearman) and \hydnIsoboleWells{} isobole wells (the growth
front). A layout; the counts are the trainer's. \status{todo}\,\fstat{idea}
\textbf{d}, The model in three stages, its equations the trainer's
(\texttt{MixtureModel} of \file{scripts/train_mixture.py}), set by LaTeX with each
component named under its brace. One inhibitor $m_i$ (furfural, drawn as in a) is
featurized as its \hydnFingerprintBits{} FCFP4 counts and passed through the compound
MLP to one token $t_i$. The token is placed in front of the \hydnGeneTokens{} gene
tokens $G$, a learned table read through the one-layer cell graph transformer and
trained with the rest at its own rate, not the composite sequence embedding of the
interactions figure; one further transformer layer reads the sequence, so the token
reads the genome and the genes read the compound, giving the cell-in-medium state $H^c$
and the read token $t_i'$. One head reads a state made of the rows of $H^c$ at the
strain's deleted genes $D$, the mean over genes and $t_i'$ (and, not drawn, the
encoder's own post-deletion rows and a learned source token), scaled and shifted per
source; a strain gives $\hat r_{g,i}$, its response, which over strains $g$ is the
profile the centered Spearman ranks; the host, which deletes nothing, gives $\hat f_i$,
its growth. There is no cell token: the host's state is the gene mean. The dose enters
once, as a gain $1 + g(d_i)$ on the medium offset $o(t_i)$, a per-medium bias added to
the head, with $g$ an MLP of the standardized log molar dose started at zero (the dose
offset of round ten). A mixture is a set of tokens, each read alone through that head,
and the readings compose as Bliss independence: a product on the host's fitness scale, a
sum on a screen's log response. The model is then scored on the \hydnCombinations{}
combinations it never saw. A layout; the counts are the trainer's.
\status{todo}\,\fstat{idea}
\textbf{e}, The ladder, ten best arms of \hydladderNArms{} on Vanacloig's compound-cold
folds: nested ridge on fingerprint counts reads \hydladderRidge{}, the next arm
\hydladderSecond{}, the tenth \hydladderTenth{}. \status{todo}\,\fstat{file}
\textbf{f}, Which model predicts an unseen Vanacloig compound's profile. The second
round is complete: over \hydbilinearEvaluations{} compound-evaluations,
\hydenvencCompounds{} compounds on each of three fold seeds, the three-seed ensemble of
the bilinear control reads \hydbilinearMedian{} and of the environment encoder
\hydenvencMedian{}, both below nested ridge at \hydladderRidge{}. Paired on compound and
fold, the encoder minus its control is \hydpairedMedian{} with \hydpairedUp{} of
\hydpairedEvaluations{} up. The last bar is the mixture trainer's single-source arm on
its own basis (28 compounds with a molar dose, folds 0 to 2, the three-seed ensemble):
\hydmixVanEnsemble{}. \status{todo}\,\fstat{file}
\textbf{g}, The learning curve in fitted compounds. Ridge: \hydcurveEight{} at eight to
\hydcurvePool{} on the pool, about \hydcurvePerDoubling{} per doubling. The environment
encoder (038 round 17, \hydcurveEncEvaluations{} compound-evaluations per point):
\hydcurveEncEight{}, \hydcurveEncOneSix{} and \hydcurveEncTwoFour{} at eight, sixteen
and twenty-four, below ridge at every count and closing on it.
\status{todo}\,\fstat{file}
\textbf{h}, Dose, the host's potency against the screen's: furfural
\hydicThirtyFf{} millimolar against \hydvanacloigFf{}, ratio \hydratioFf{}; 5-HMF
\hydicThirtyHmf{} against \hydvanacloigHmf{}, ratio \hydratioHmf{}; the other four are in
no screen. \status{todo}\,\fstat{file}
\textbf{i}, How a pair's profile composes from its singles over \hydhetPairs{}
Hillenmeyer HET pairs: the mean is best in \hydhetMeanWins{}, median R-squared
\hydhetrTwomean{} against \hydhetrTwosum{} for the sum and \hydhetrTwomax{} for the
highest single. \status{todo}\,\fstat{file}
\textbf{j}, The predicted profile against the measured one. Furfural: the screen repeats
itself at \hydprofFfReliability{}, ridge reads \hydprofFfRidge{}; 5-HMF:
\hydprofHmfReliability{}, which is noise, against \hydprofHmfRidge{}. The two profiles
share \hydgoSharedTermsStore{} significant terms. The cell graph transformer, every
round-3 and round-4 run on both machines in which the compound was held out: furfural
\hydprofFfCgt{} over \hydprofFfCgtRuns{} runs, 5-HMF \hydprofHmfCgt{} over
\hydprofHmfCgtRuns{}, with the run spread as the whisker. \status{todo}\,\fstat{file}
\textbf{k}, The pre-registered test, which FAILS: does the similarity of two
inhibitors' profiles anticipate how far their pair departs from Bliss. With measured
profiles where one exists the association is \hydsvdBestSpearman{} at a permutation p
of \hydsvdBestP{} over \hydsvdBestPairs{} pairs; with every profile predicted by ridge
it is \hydsvdPredSpearman{} at p \hydsvdPredP{}. \status{todo}\,\fstat{file}
\textbf{l}, This study's six titrations with their Hill fits, dose in units of each
compound's own half-maximal dose; slopes run \hydhillMin{} to \hydhillMax{}. The model
whose host, the empty genotype read through the shared cell head, is trained on the
public anchors alone reads a flat line at the anchors' assumed fitness for every
inhibitor at every dose, Spearman \hydexTwoOneZeroShotSpearman{} against the
\hydexTwoOneZeroShotConditions{} observed single-agent conditions (not drawn); trained
on these wells it fits, Spearman \hydexTwoOneFittedSpearman{}, and is still nearly flat
in dose, because without a dose channel the host reads one value per compound (not
drawn). Drawn, the same model with the dose offset (round ten): the compound tokens
stay dose-free and the dose scales the medium offset, so the curve falls with dose,
Spearman \hydexTwoOneDoseSpearman{} on these wells. \status{todo}\,\fstat{file}
\textbf{m}, This study's combinations that grew, observed against the two rules, by
number of inhibitors. Growth AUROC: Loewe \hydloeweAuroc{}, Bliss \hydblissAuroc{}.
Fitness Spearman: Loewe \hydloeweSpearman{}, Bliss \hydblissSpearman{}. The fourth bar
is the model (single source, cell readout, trained on the anchors and the titration,
mean over \hydmodelRuns{} runs on \hydmodelFolds{} folds, seeds and machines), which
reads a medium one inhibitor at a time through the shared head and composes the
readings as Bliss independence: its predicted fitness falls with the inhibitor count,
growth AUROC \hydmodelAuroc{}, fitness Spearman \hydmodelGrownSpearman{}. The fifth bar,
lighter, is the same model with the dose offset (\hydmodelDoseRuns{} runs): it reads
the titration (l) and still composes, growth AUROC \hydmodelDoseAuroc{}, fitness
Spearman \hydmodelDoseGrownSpearman{}, with the fitness at one inhibitor at the observed
level and a shallower fall at two and three. Under mean pooling of the inhibitor tokens
the same model read one fitness at every count (AUROC 0.535).
\status{todo}\,\fstat{file}
\textbf{n}, This study's \hydnIsoboles{} isobole grids as observed fitness, white at
none, each grid's inhibitor named on its lower edge and acetic acid on every grid's
left edge, with the cells that grew marked and each rule's front drawn where its
predicted fitness crosses the growth threshold: Bliss in red, Loewe in purple, every
front a solid line over a white halo. Mean observed minus Bliss over the
\hydisoboleExTwoSixCells{} interior cells: \hydisoboleExTwoSixExcess{},
\hydisoboleExTwoSevenExcess{} and \hydisoboleExTwoEightExcess{}, the last a single run
with growth islands and flagged. In blue, the Bliss-composed model's front after the
monotone closure, inside the grid for the first time: the model's mean observed minus
predicted over the interior cells is -0.084, +0.016 and -0.102, against the rule's
-0.082, -0.446 and -0.017, so on the formic acid grid, where the rule on measured
singles overpredicts by 0.45, the model is within 0.02. The dose-offset arm is scored on
the same grids (supplement eight) and not drawn. A caveat the sheet carries: the
model's single-agent curves are nearly flat in dose (l), so a product of two low
readings matching a low interior is not yet evidence that it reads the interaction; the
closure drops \hydisoboleExTwoSixModelIslands{} cells on the furfural grid.
\status{todo}\,\fstat{file}
\textbf{o}, The public screens added to the task under the Bliss cell readout, every
fitted and zero-shot arm on the same rows, the mean over every run on fold seed 0 with
the run spread as the whisker. Left, the compound-cold centered Spearman: one source
\hydrOneZeroonegenemean{} (\hydrOneZeroonegenen{} runs), with Hoepfner HOP
\hydrOneZerohopgenemean{}, with Hillenmeyer HET \hydrOneZerohetgenemean{}, with
Wildenhain \hydrOneZerowilgenemean{}, three \hydrOneZerothreegenemean{}, four
\hydrOneZerofourgenemean{}, every one inside one spread and below nested ridge at
\hydladderRidge{}; no screen lifts the task. Right, the growth AUROC of the same arms
over the \hydnCombinations{} combinations. The fitted arms, which train on the
titration, read \hydrOneZeroonegrowthmean{} with one source to
\hydrOneZerofourgrowthmean{} with four. The zero-shot arms, which never see a titration
well, read \hydrOneZerozsOnegrowthmean{}, \hydrOneZerozsTwogrowthmean{} and
\hydrOneZerozsThreegrowthmean{} on one, two and three sources
(\hydrOneZerozsOnegrowthn{} runs each), against Loewe at \hydloeweAuroc{}: a rise with
every screen added, with the steps inside the spread, so a measured ordering rather
than a cleared difference. \status{todo}\,\fstat{file}

%% The float is flushed before its source table is set: a deferred float
%% otherwise lands on the page AFTER the table that describes it.
\clearpage

{\footnotesize\noindent
\begin{tabular}{@{}p{8mm}p{168mm}@{}}
\toprule
\multicolumn{2}{@{}l}{\textbf{Sources for \cref{fig:main}}} \\
\midrule
a & \file{profile_fingerprint_check.csv} (the curated SMILES of the pool's compounds;
the sheet computes each inhibitor's nearest by FCFP4 count Tanimoto, itself and the
other inhibitors excluded) and \file{source_inhibitor_overlap.csv} (which inhibitors a
screen profiles), and
\file{scripts/mixture_data.py} \texttt{IDENTITY\_GAP\_SMILES} for formic and lactic
acid; \file{single_agent_fits.csv} for the half-maximal doses. The structures are
RDKit's ACS 1996 drawings as vector paths (\texttt{torchcell.utils.acs\_structure}).
Status: on file. \\
b & \file{results/mixture-delta/r3_vanhophetwil_e150_p0/*_f0_counts.csv}, the
per-source counts \file{scripts/train_mixture.py} writes beside every run;
\file{source_label_histograms.csv} and \file{source_label_summary.csv}
(\file{scripts/source_labels.py}, the oriented value per source read from the 033 cell
table and from \file{wetlab_wells.csv}); \file{wetlab_wells.csv},
\file{single_agent_fits.csv}, \file{mixture_combinations.csv} and
\file{isobole_cells.csv} for this study's rows drawn as they are. The genotypes and the
culture conditions are the loaders' sourced values (\file{vanacloig2022.py}, \file{hoepfner2014.py}, \file{hillenmeyer2008.py},
\file{wildenhain2015.py}, \file{private_torchcell/bioscreen.py}; the Bioscreen archive
states no shaking, so ``shaken'' is the operator's statement). ``This study'' is the
private Bioscreen archive: the titration is run ex21, the combinations ex23, the
isobole grids ex26 to ex28. Status: on file. \\
c & the same counts file; \file{source_agreement.csv} (\file{scripts/source_axes.py},
the compounds two screens share, from 033's cross-source cells);
\file{source_inhibitor_overlap.csv} (\file{scripts/source_labels.py}, the inhibitors a
screen holds, matched on InChIKey); \file{scripts/mixture_data.py}
\texttt{anchor\_records} for the anchors. The drawing is a layout. Status: idea. \\
d & \file{scripts/train_mixture.py} (the architecture, \texttt{MixtureModel}) and
\file{scripts/mixture_data.py} (the fingerprint); the drawing is a layout. Status:
idea. \\
e & \file{038/ladder/ladder_r2_summary.csv}, centered target, full compound count.
Status: on file. \\
f & \file{038/ladder/ladder_r2_summary.csv}; \file{038/factorized/r9_control/*_scores.csv};
\file{038/factorized/r10_envenc/*_scores.csv}; \file{results/mixture/r3_van_e150/*_scores.csv}
for the mixture-trainer bar. Status: on file. \\
g & \file{035/learning_curve/ridge_curve.csv}, centered target;
\file{038/factorized/r17_curve/*_scores.csv} for the encoder. Status: on file. \\
h & \file{single_agent_fits.csv} (the host, with bootstrap intervals from
\file{single_agent_bootstrap.csv}); \file{vanacloig_ic30.csv} (the screen). Status: on
file. \\
i & \file{het_rule_fit.csv}, \file{het_summary.json}. Status: on file. \\
j & \file{profile_predicted_vs_measured.csv} (\texttt{loco\_32}, centered);
\file{go_overlap_measured_vs_predicted.csv}; every
\file{results/mixture*/r[34]*/*_vanacloig_compounds.csv} on both machines for the
transformer. Status: on file. \\
k & \file{similarity_vs_deviation_pairs.csv}, \file{similarity_vs_deviation.csv}
(both profile versions, the served growth call). Status: on file. \\
l & \file{wetlab_wells.csv} (run ex21), \file{single_agent_fits.csv},
\file{single_agent_bootstrap.csv}; \file{results/host_predictions.csv}
(\file{scripts/r3_host_predictions.py}) for the model curves. Status: on file. \\
m & \file{mixture_combinations.csv}, \file{mixture_scores.csv} (run ex23, served call);
\file{results/host_predictions.csv} for the model bars. Status: on file. \\
n & \file{isobole_cells.csv}, \file{isobole_summary.csv} (runs ex26 to ex28);
\file{results/host_predictions.csv} and \file{mixture_rules.monotone_alive} for the
model front. Status: on file. \\
o & every \file{results/mixture*/r10_*_cell_bliss*/*_scores.csv} on both machines
(fold seed 0): the \texttt{centered\_spearman\_median} rows on the held-out compounds
(left) and the \texttt{growth\_auroc} rows on the served call of the ex23 combinations
(right), with the anchors-only arms beside the fitted ones. Status: on file. \\
\bottomrule
\end{tabular}}

\clearpage
\tcfig{\figSiOne}{\textbf{Supplement one: the public benchmark in full.}
\textbf{a}, Every ladder arm, median centered Spearman per encoder and kernel,
\hydladderCells{} fitted cells over \hydladderEncoders{} encoders. Fingerprint counts
lead every kernel they are defined for. \status{todo}\,\fstat{file}
\textbf{b}, Each of the \hydperCompoundN{} compounds, its ridge score against its own
reliability ceiling: median \hydperCompoundMedian{}, \hydperCompoundAbove{} above the
ceiling and \hydperCompoundNegative{} below zero. \status{todo}\,\fstat{file}
\textbf{c}, The learning curve on both targets, the uncentered one dashed. The centered
curve is the one the main sheet carries. \status{todo}\,\fstat{file}
\textbf{d}, What each rule gets wrong, by combination size: how many combinations the
rule called alive, how many of those were right, and what it missed or invented. Loewe
misses all \hydloeweMissedThree{} growing triples and \hydloeweMissedTwo{} of
\hydloeweGrewTwo{} growing pairs; Bliss calls \hydblissFalseThree{} of
\hydblissDeadThree{} dead triples alive. \status{todo}\,\fstat{file}}{fig:sione}

%% The float is flushed before its source table is set: a deferred float
%% otherwise lands on the page AFTER the table that describes it.
\clearpage

{\footnotesize\noindent
\begin{tabular}{@{}p{8mm}p{168mm}@{}}
\toprule
\multicolumn{2}{@{}l}{\textbf{Sources for \cref{fig:sione}}} \\
\midrule
a & \file{038/ladder/ladder_r2_summary.csv}. Status: on file. \\
b & \file{038/ladder/ladder_r2_scores.csv}, kernel ridge on fingerprint counts,
centered, averaged over fold seeds. Status: on file. \\
c & \file{035/learning_curve/ridge_curve.csv}. Status: on file. \\
d & \file{mixture_scores.csv}, served call, one row per combination size. Status: on
file. \\
\bottomrule
\end{tabular}}

\clearpage
\tcfig{\figSiTwo}{\textbf{Supplement two: the single agents and the growth
call.}
\textbf{a}, Every single-agent fit, \hydnFitsAll{} of them over the titration run and the
three isobole axes, under both no-growth conventions: the darker bar keeps no-growth
wells at zero fitness, the lighter one drops them from the fit. Intervals are the
bootstrap. \status{todo}\,\fstat{file}
\textbf{b}, Every claim-one score under both growth calls, hatched for the software call.
The two calls disagree on \hydnDisputedWells{} wells of the combination run across
\hydnDisputedCombinations{} combinations, with optical-density rises from
\hydodRiseMin{} to \hydodRiseMax{}. \status{todo}\,\fstat{file}
\textbf{c}, Each pair against Bliss with its replicate-bootstrap interval, under both
calls. Under the served call \hydnPairsSynergy{} of \hydnPairs{} pairs are synergistic,
\hydnPairsAdditive{} additive and \hydnPairsAntagonism{} antagonistic.
\status{todo}\,\fstat{file}
\textbf{d}, Where the Loewe rule had to extrapolate past the measured dose range:
\hydloeweBeyondTotal{} of the combinations, and \hydisoboleFrontsBeyondGrid{} of the
isobole fronts fell outside their grid. \status{todo}\,\fstat{file}}{fig:sitwo}

%% The float is flushed before its source table is set: a deferred float
%% otherwise lands on the page AFTER the table that describes it.
\clearpage

{\footnotesize\noindent
\begin{tabular}{@{}p{8mm}p{168mm}@{}}
\toprule
\multicolumn{2}{@{}l}{\textbf{Sources for \cref{fig:sitwo}}} \\
\midrule
a & \file{single_agent_fits.csv}, \file{single_agent_bootstrap.csv}. Status: on
file. \\
b & \file{mixture_scores.csv} (both calls), \file{wetlab_growth_call_check.csv}.
Status: on file. \\
c & \file{pair_deviation.csv}. Status: on file. \\
d & \file{mixture_combinations.csv} (the extrapolation flags),
\file{isobole_summary.csv}. Status: on file. \\
\bottomrule
\end{tabular}}

\clearpage
\tcfig{\figSiThree}{\textbf{Supplement three: how profiles compose, and what the
predicted ones carry.}
\textbf{a}, Emergent and masked gene rates of each public pair against how alike its two
singles are, with the control-matched null beside them. The median emergent fraction is
\hydhetEmergentMedian{} against a null of \hydhetNullEmergent{}, and the masked fraction
\hydhetMaskedMedian{}; the only two near-replicate singles in the screen disagree on
\hydhetReplicateDisagreeMin{} to \hydhetReplicateDisagreeMax{} of their own hits, which
is the noise floor these rates sit inside. \status{todo}\,\fstat{file}
\textbf{b}, The one dose grid the public screen holds, methotrexate against
fluorouracil: the weight the free linear fit puts on the first single moves from
\hydmtxLinAMin{} to \hydmtxLinAMax{} across \hydmtxCells{} dose pairs, while the mean
wins at \hydmtxMeanWins{} of them. \status{todo}\,\fstat{file}
\textbf{c}, \textbf{d}, Profile similarity among the six inhibitors, measured where a
measurement exists and then with every profile predicted. Measured, the strongest pair
reads \hydsimBestAvailableMax{}; predicted, \hydsimAllPredictedMax{}, because the small
acids have near-identical predicted profiles. One scale carries both panels.
\status{todo}\,\fstat{file}
\textbf{e}, Similarity against the pair's departure from Bliss, both versions:
\hydsvdBestSpearman{} at a permutation p of \hydsvdBestP{} with measurements in, and
\hydsvdPredSpearman{} at \hydsvdPredP{} with everything predicted.
\status{todo}\,\fstat{file}
\textbf{f}, The GO read: significant terms in the top hundred sensitive genes of the
measured profile, of the ridge-predicted profile, and of the gene mean over the fitted
compounds. The control carries \hydgoControlTerms{} terms and the predicted profiles
share \hydgoSharedControlMax{} of them, which is what the predicted enrichment is.
\status{todo}\,\fstat{file}}{fig:sithree}

%% The float is flushed before its source table is set: a deferred float
%% otherwise lands on the page AFTER the table that describes it.
\clearpage

{\footnotesize\noindent
\begin{tabular}{@{}p{8mm}p{168mm}@{}}
\toprule
\multicolumn{2}{@{}l}{\textbf{Sources for \cref{fig:sithree}}} \\
\midrule
a & \file{het_emergent_masked.csv}, \file{het_near_replicates.csv},
\file{het_summary.json}. Status: on file. \\
b & \file{het_rule_fit.csv}, the methotrexate and fluorouracil family. Status: on
file. \\
c, d & \file{profile_similarity.csv}, centered target, both versions. Status: on file. \\
e & \file{similarity_vs_deviation_pairs.csv}, \file{similarity_vs_deviation.csv}.
Status: on file. \\
f & \file{go_overlap_measured_vs_predicted.csv}, \file{nominations_summary.csv}.
Status: on file. \\
\bottomrule
\end{tabular}}

\clearpage
\tcfig{\figSiFour}{\textbf{Supplement four: what the mixture rounds say, read
over both machines.} Rounds three and four of the mixture trainer, every run on fold
seed 0 unless a panel says otherwise; a run is one seed on one fold on one machine, and
an arm run on both machines contributes both copies.
\textbf{a}, The one-, two- and three-source arms on the public side, mean over runs with
its standard deviation, \hydrThreeonegenen{}, \hydrThreetwogenen{} and
\hydrThreethreegenen{} runs: \hydrThreeonegenemean{}, \hydrThreetwogenemean{} and
\hydrThreethreegenemean{}, against nested ridge at \hydladderRidge{}. The run spread is
\hydrThreeonegenesd{} to \hydrThreethreegenesd{}. \status{todo}\,\fstat{file}
\textbf{b}, The same runs on the wet lab, growth AUROC over the 63 combinations under
the served call: \hydrThreeonegrowthmean{}, \hydrThreetwogrowthmean{} and
\hydrThreethreegrowthmean{}, against Loewe at \hydloeweAuroc{}. These are the gene
readout's arms; the same ladder under the Bliss cell readout, at three seeds, is
\cref{fig:sieight}. \status{todo}\,\fstat{file}
\textbf{c}, The dose channel, paired: the same arm with the log-molar dose modulation on
and off, same machine, same folds, one seed; \hydrThreedosepairs{} pairs, FiLM lower on
\hydrThreedosepairsfilmLower{} of them, and the furfural-grid residual more negative on
\hydrThreedosepairsgridMoreNegative{} of them. \status{todo}\,\fstat{file}
\textbf{d}, What the transformer predicts on the public screens it learns from: the
held-out compound score per source, mean over the plain arms that carry it, beside the
chemistry-feature learner's compound-cold score on the same source (Hoepfner
\hydrThreesrchoepfnercgt{} against \hydrThreesrchoepfnerridge{}, Hillenmeyer
\hydrThreesrchillenmeyercgt{} against \hydrThreesrchillenmeyerridge{}, Wildenhain
\hydrThreesrcwildenhaincgt{} against \hydrThreesrcwildenhainridge{}); and the held-out
DOSE of a seen Hoepfner compound under FiLM, \hydrThreesrchoepfnerdose{}.
\status{todo}\,\fstat{file}
\textbf{e}, The same config at the same seed on the two machines, every round,
\hydrThreemachinesn{} config-fold pairs: median absolute difference
\hydrThreemachinesmedianAbsDiff{}, largest \hydrThreemachinesmaxAbsDiff{}, Spearman
between the copies \hydrThreemachinesspearman{}. That is the run-to-run floor at a
fixed seed. \status{todo}\,\fstat{file}
\textbf{f}, The partition: each source-set arm on three fold seeds, one seed, median
over the folds on file. The single-source arm reads \hydrThreefsoneZero{},
\hydrThreefsoneOne{} and \hydrThreefsoneTwo{}. \status{todo}\,\fstat{file}}{fig:sifour}

%% The float is flushed before its source table is set: a deferred float
%% otherwise lands on the page AFTER the table that describes it.
\clearpage

{\footnotesize\noindent
\begin{tabular}{@{}p{8mm}p{168mm}@{}}
\toprule
\multicolumn{2}{@{}l}{\textbf{Sources for \cref{fig:sifour}}} \\
\midrule
a, b & every \file{results/mixture/r3*/*_scores.csv} and \file{r4*/*_scores.csv} on
GilaHyper and \file{results/mixture-delta/r3*/*_scores.csv} pulled from Delta, seed
members of the arms \file{mix_van_e150}, \file{mix_vanhop_e150} and
\file{mix_vanhophet_e150}; the ridge bar from the 038 ladder, the Loewe bar from
\file{mixture_scores.csv}. Status: on file. \\
c & the same files, the \file{_film} arm against its dose-off twin on the folds both have,
seed 0. Status: on file. \\
d & the same files, rows whose task is the auxiliary source (\file{heldout_compound},
\file{heldout_dose}); the chemistry-feature bars from \file{source_basis.csv}. Status: on
file. \\
e & the same files, every config and fold scored on both machines at seed 0. Status: on
file. \\
f & the same files, the \file{_fs1} and \file{_fs2} arms. Status: on file. \\
\bottomrule
\end{tabular}}

\clearpage
\tcfig{\figSiFive}{\textbf{Supplement five: is the trainer converged, per
source.} Mean curves over every round-3 and round-4 run on both machines of the
one-, two- and three-source arms and two three-source variants; the learning rate is zero
at epoch 150 on every arm. \textbf{a}, The training loss per source, one step per
evaluation, smoothed. The target's falls \hydconvVanTrainEarly{} to
\hydconvVanTrainEnd{}; Hoepfner's \hydconvHopTrainEarly{} to \hydconvHopTrainEnd{};
Hillenmeyer's RISES, \hydconvHetTrainEarly{} to \hydconvHetTrainEnd{}.
\status{todo}\,\fstat{file}
\textbf{b}, The held-out loss per source. The target's removes \hydconvVanRemoved{} of its
first value, bottoms near epoch \hydconvVanMinEpoch{} and rises \hydconvVanRise{} by the
end; Hoepfner's removes \hydconvHopRemoved{}; Hillenmeyer's \hydconvHetRemoved{};
Wildenhain's \hydconvWilRemoved{}, bottoming near epoch \hydconvWilMinEpoch{} and still
falling at 150. \status{todo}\,\fstat{file}
\textbf{c}, The held-out and validation centered Spearman by epoch: the held-out score
plateaus near epoch 60 to 75 while the held-out loss has already turned, so the ranking
holds as the scale drifts. \status{todo}\,\fstat{file}}{fig:sifive}

\clearpage

{\footnotesize\noindent
\begin{tabular}{@{}p{8mm}p{168mm}@{}}
\toprule
\multicolumn{2}{@{}l}{\textbf{Sources for \cref{fig:sifive}}} \\
\midrule
a to c & every \file{results/mixture/r3*/*_history.csv} and \file{r4*} on GilaHyper and
\file{results/mixture-delta/r3*/*_history.csv} from Delta, read by
\file{scripts/r3_convergence.py} into \file{results/r3_convergence_runs.csv} and
\file{r3_convergence_arms.csv}. Status: on file. \\
\bottomrule
\end{tabular}}

\clearpage
\tcfig{\figSiSix}{\textbf{Supplement six: every scored arm of rounds one to
four, both machines.} One row per arm, every run a point (one seed on one fold on one
machine) and the arm's mean a bar, sorted by the public score; rounds one and two
(pool-and-last, 50 epochs) sit on the same axis as rounds three and four (the held-out
protocol, 150 epochs). \textbf{a}, The compound-cold centered Spearman on Vanacloig
against nested ridge at \hydladderRidge{}. \textbf{b}, The growth call over the 63
combinations against Loewe at \hydloeweAuroc{}. The spread of the points within an arm
against the range of the arm means is the noise floor: no arm clears ridge and none
approaches Loewe. \status{todo}\,\fstat{file}}{fig:sisix}

\clearpage

{\footnotesize\noindent
\begin{tabular}{@{}p{8mm}p{168mm}@{}}
\toprule
\multicolumn{2}{@{}l}{\textbf{Sources for \cref{fig:sisix}}} \\
\midrule
a, b & every \file{*_scores.csv} under \file{results/mixture/} and
\file{results/mixture-delta/}, drawn by \file{scripts/r3_plots.py}; the reference values
from \file{results/figure_gate_numbers.json}. Status: on file. \\
\bottomrule
\end{tabular}}

\clearpage
\subsection{The encoders, and how they were compared}
\label{sec:encoders}

Every encoder in this document was scored on ONE basis, so the sheet's bars are
comparable: a compound-cold split of the Vanacloig profiles (the compound is held out
with every one of its doses; 32 compounds on each of three fold seeds, 96
compound-evaluations), the centered Spearman across strains between the predicted and
the measured profile of the held-out compound, and the median over the evaluations as
the score. Three kinds of encoder sit on that basis.

\begin{itemize}
\item \textbf{Fixed chemistry encoders with a shallow learner} (the ladder, 038):
  \hydencoderN{} ways of embedding a compound, four fingerprint families (FCFP4 and
  ECFP4 as counts or bits, MACCS keys), RDKit's 2D descriptors, seven pretrained
  molecular models (mol2vec, MolE, Uni-Mol, ChemBERTa-2 in both pretraining modes,
  MoLFormer-XL, a ZINC RoBERTa) and three concatenations, each run with kernel ridge
  and nearest neighbors under linear, RBF and Tanimoto kernels, \hydencoderLearnerRuns{}
  learner-by-encoder arms, every regularization chosen by a nested fold. The best is
  nested ridge on FCFP4 counts at \hydencoderBest{}; the best pretrained model is
  \hydencoderBestPretrainedName{} at \hydencoderBestPretrained{}; the lowest single
  encoder is \hydencoderWorstName{} at \hydencoderWorst{}.
\item \textbf{Learned environment encoders} (038 rounds 9 and 10): the bilinear control
  (a compound embedding read against a learned gene embedding, no attention) and the
  environment encoder (the compound token with its dose in front of the gene tokens of
  the cell graph transformer), trained on the same folds, three seeds each; the
  encoder is read against its control paired on compound and fold, median difference
  \hydpairedMedian{}, up on \hydpairedUp{} of \hydpairedEvaluations{}.
\item \textbf{The mixture trainer} (040), the same transformer with the public screens
  and the host records in one objective, which scores on its own basis (the 28
  compounds with a molar dose, folds 0 to 2): \hydmixVanEnsemble{} for the gene readout
  and, under the cell readout, inside its fold noise (\cref{sec:state}).
\end{itemize}

\tcfig{\figSiSeven}{\textbf{Supplement seven: the encoders compared.} \textbf{a}, Every
chemistry encoder of the ladder on the one basis, its best learner as the bar (the
learner named in the label) and every learner it was run with as a point; color is the
encoder family: fingerprints, descriptors, pretrained models, concatenations. Nested
ridge on fingerprint counts leads; no pretrained model reaches it. \textbf{b}, The
learned environment encoder against its bilinear control, paired on compound and fold
over the \hydpairedEvaluations{} evaluations, sorted by the control's score; the line
is nested ridge. The encoder is not above its control (median difference
\hydpairedMedian{}) and neither is above ridge. \status{todo}\,\fstat{file}}{fig:siseven}

\clearpage

{\footnotesize\noindent
\begin{tabular}{@{}p{8mm}p{168mm}@{}}
\toprule
\multicolumn{2}{@{}l}{\textbf{Sources for \cref{fig:siseven}}} \\
\midrule
a & \file{038/ladder/ladder_r2_summary.csv}, the centered target at the full compound
count, every learner-by-encoder row. Status: on file. \\
b & \file{038/factorized/r9_control/*_scores.csv} and
\file{038/factorized/r10_envenc/*_scores.csv}, the three-seed ensemble per fold, merged
on compound, fold and fold seed. Status: on file. \\
\bottomrule
\end{tabular}}

\clearpage
\tcfig{\figSiEight}{\textbf{Supplement eight: round ten, the Bliss cell readout at
three seeds.} Every arm is three seeds on the three screening folds of fold seed 0,
read on both machines where it ran on both; the bar is the mean over runs, the whisker
the standard deviation.
\textbf{a}, The source ladder on the public side: one source \hydrOneZeroonegenemean{},
with Hoepfner \hydrOneZerohopgenemean{}, with Hillenmeyer \hydrOneZerohetgenemean{}, with
Wildenhain \hydrOneZerowilgenemean{}, three \hydrOneZerothreegenemean{}, four
\hydrOneZerofourgenemean{}, run spreads \hydrOneZeroonegenesd{} to \hydrOneZerofourgenesd{}, against
nested ridge at \hydladderRidge{}. No public screen lifts the target, and the order is
more sources, lower, inside one spread. \status{todo}\,\fstat{file}
\textbf{b}, The same arms on the wet lab, growth AUROC over the 63 combinations: the
fitted ladder reads \hydrOneZeroonegrowthmean{} to \hydrOneZerofourgrowthmean{} and does not
move with sources; the zero-shot ladder (anchors only, never a titration well) reads
\hydrOneZerozsOnegrowthmean{}, \hydrOneZerozsTwogrowthmean{} and \hydrOneZerozsThreegrowthmean{} on
one, two and three sources, every value above the Bliss rule on the measured singles
(\hydblissAuroc{}) and the last at Loewe's (\hydloeweAuroc{}). Three runs per arm; the
steps are inside one spread, so this is a measured ordering rather than a cleared
difference. \status{todo}\,\fstat{file}
\textbf{c}, The dose channel: each arm's growth AUROC (dark) beside the Spearman of its
host prediction against the 180 titration wells (light), a fit for the arms trained on
the titration and a prediction for the zero-shot arms. FiLM on the tokens reads the
titration at \hydrOneZerodosefilmtitration{} and composes at \hydrOneZerodosefilmauroc{}; the
dose offset reads it at \hydrOneZerodosedofftitration{} and composes at
\hydrOneZerodosedoffauroc{}, against \hydrOneZerodoseoffauroc{} with no channel; the titration
fine-tune with FiLM reads \hydrOneZerodosefttitration{} and composes at
\hydrOneZerodoseftauroc{}. Zero-shot, no channel reads the titration, and the dose offset
composes at \hydrOneZerodosezsDoffauroc{}. \status{todo}\,\fstat{file}
\textbf{d}, What the Bliss arms predict on the public screens they learn from, the
held-out compound score per source over every round-ten arm that carries it: Hoepfner
\hydrOneZerosrchoepfnercgt{} against \hydrOneZerosrchoepfnerridge{}, Hillenmeyer
\hydrOneZerosrchillenmeyercgt{} against \hydrOneZerosrchillenmeyerridge{}, Wildenhain
\hydrOneZerosrcwildenhaincgt{} against \hydrOneZerosrcwildenhainridge{}; the held-out Hoepfner
dose under FiLM, \hydrOneZerosrchoepfnerdose{}. The screens are not learned at a held-out
compound, so a screen could not have helped the target at one.
\status{todo}\,\fstat{file}}{fig:sieight}

\clearpage

{\footnotesize\noindent
\begin{tabular}{@{}p{8mm}p{168mm}@{}}
\toprule
\multicolumn{2}{@{}l}{\textbf{Sources for \cref{fig:sieight}}} \\
\midrule
a, b & \file{results/mixture-delta/r10*/*_scores.csv} pulled from Delta and
\file{results/mixture/r9*/*_scores.csv} on GilaHyper (the seed-0 copies of the one-
and three-source Bliss arms), seed members on fold seed 0; the ridge bar from the 038
ladder, the Loewe bar from \file{mixture_scores.csv}. Status: on file. \\
c & the same files, the \file{ex23} growth AUROC and the \file{ex21} all-wells fitness
Spearman of the seven dose arms. Status: on file. \\
d & the same files, rows whose task is the auxiliary source (\file{heldout_compound};
\file{heldout_dose} of \file{mix_vanhop_e150_cell_bliss_film_dose}); the
chemistry-feature bars from \file{source_basis.csv}. Status: on file. \\
\bottomrule
\end{tabular}}
